\qquad Este trabalho tem como objetivo implementar e analisar quatro
algoritmos clássicos de ordenação por comparação estudados na
disciplina de Projeto de Algoritmos: Insertion Sort, Selection Sort,
Bubble Sort e Shell Sort. Foi desenvolvida uma aplicação em C++,
com criação automática das pastas e arquivos pelo próprio programa,
capaz de gerar instâncias de teste de diferentes tamanhos (10, 100,
1.000, 10.000, 100.000 e 1.000.000 de elementos) e diferentes
características (crescente, decrescente e randômica), executar cada
algoritmo sobre elas e registrar o tempo de execução observado.

\qquad Os resultados foram organizados em arquivos de entrada, saída e
tempo para cada algoritmo, consolidados em tabelas e gráficos
individuais e comparativos que permitem analisar o comportamento de
cada algoritmo isoladamente e compará-los entre si. As medições
confirmam o comportamento teórico esperado para cada algoritmo:
desempenho linear no melhor caso do Insertion, Bubble e Shell Sort
(entrada já ordenada), desempenho quadrático constante em qualquer
caso para o Selection Sort, e uma clara superioridade prática do
Shell Sort sobre os demais algoritmos $O(n^2)$, evidenciando a forte
influência da estratégia de cada algoritmo — e, no caso do Insertion,
Selection e Bubble Sort, também do grau de ordenação inicial dos
dados — sobre o desempenho observado.
